\documentclass[addpoints,ngerman,12pt,answers]{exam}
\usepackage{babel}
\usepackage[top=2cm,bottom=2.5cm,left=1.5cm,right=1.5cm]{geometry}
\usepackage[utf8]{inputenc}
\usepackage[T1]{fontenc}
\usepackage{booktabs}
\usepackage{graphicx}
\usepackage{csquotes}
\usepackage{paralist}
%\usepackage{wasysym}
%\usepackage[math]{iwona}
\usepackage{setspace}
\usepackage{amsmath} 
\usepackage[math]{iwona}
\usepackage{setspace}
\usepackage{enumitem}
\usepackage{amsmath}
\usepackage{amsfonts}


\pointpoints{Point}{Points}
\bonuspointpoints{Bonuspunkt}{Bonuspunkte}
\renewcommand{\solutiontitle}{\noindent\textbf{Solution:}\enspace}

\chqword{Frage}   
\chpgword{Seite} 
\chpword{Punkte}   
\chbpword{Bonus Points} 
\chsword{Erreicht}   
\chtword{Gesamt}

\checkboxchar{\Square}
\checkedchar{\CheckedBox}

\pagestyle{headandfoot}
\runningheadrule
\firstpageheader{ECON 311}{Tutorial 1 and 2}{}
\runningheader{ECON 311}{Tutorial 1 and 2}{}
\firstpagefooter{ECON 311}{Tutorial 1 and 2}{\thepage\,/\,\numpages}
\runningfooter{ECON 311}{Tutorial 1 and 2}{\thepage\,/\,\numpages}
\doublespacing

\begin{document}
	\vspace*{3em}
	
	\makebox[\textwidth]{Name:\enspace\hrulefill}
	
	\vspace*{3em}
	\begin{questions}
		
		
		\question Tom has a income of \$10,000 in the current period, and income of \$12,000 in the future period. Tom does not have to pay any tax. The real interest rate is 0.05 or 5\% per period.
		
		\begin{enumerate}[label=(\alph*)]
			\item Write down Tom's budget constraint for current period and future period. Then, derive life-time budget constraint.
			\item Tom has the following utility function: $U(c,c') = \ln(c) + 0.95\ln(c')$. Find Tom's optimal choice for consumption in current and future period. Is Tom a lender or a borrower?
			\item Now suppose Tom's current income increase to \$12,000. Find Tom's new optimal choice for consumption in current and future period. Is Tom a lender or a borrower?
			\item Draw a graph on the space of $(c,c')$ representing showing the result from part c)
		\end{enumerate}
%		\begin{solution}
%			\begin{enumerate}[label=(\alph*)]
%				\item Tom's budget constraint for current period is:
%				\[c+s = 10000\]
%				Tom's budget constraint for future period is:
%				\[c' = 12000+(1+0.05)s\]
%				Rearrange Tom's budget constraint for current period
%				\[s = 10000 - c\]
%				Plug it into the budget constraint for future period
%				\[c' = 12000+(1+0.05)(10000 - c)\]
%				\[c' + (1+0.05)c = 12000+(1+0.05)(10000)\]
%				or
%				\[\frac{c'}{(1.05)} + c = \frac{12000}{(1.05)}+10000\]
%				\item Setup the Largarian
%				\[\ln(c) + 0.95\ln(c') + \lambda\left(22500 -c' - (1.05)c\right)\]
%				Take derivative w.r.t. $c$ and $c'$
%				\[\frac{1}{c} = \lambda (1.05)\]
%				\[\frac{0.95}{c'} = \lambda \]
%				Combine the two equations:
%				\[\frac{c'}{0.95c} = 1.05 \to c' = (1.05)(0.95)c = 0.9975c\]
%				Plug it back to the lifetime budget constraint
%				\[0.9975c + (1.05)c = 22500 \to 2.0475c = 22500\]
%				\[c^* = 10989.01\]
%				and 
%				\[c'^* = 22500 - (1.05)10989.01 = 10961.54\]
%				Since $c^* = 10989.01 > 10,000$, Tom is a borrower.
%				\item Now, $c$ increase to 12000. Note that we do not have to solve the whole optimization problem again. Instead, we plug in $c' = 0.9975c$ into the new lifetime budget constraint
%				\[c' + (1.05)c = 12000+(1.05)(\mathbf{12000}) \]
%				\[0.9975c+ 1.05c = 24600 \to 2.0475c = 24600 \to c^* = 12,014\]
%				\[c'^* = 24600 - (1.05)12,014 = 11985.3\]
%				Although the income in current period has increased to the same level as in the future period, Tom is still a borrower ($c^*>12000$). It is because Tom would like to smooth out his consumption over the two periods, there is also an increase in future consumption.
%			\end{enumerate}
%		\end{solution}		
%		
		\question Textbook Question: 9.1, 9.8
	\end{questions}
	
%	\begin{center}
%		\combinedgradetable[h][questions]
%	\end{center}
	
\end{document}